\documentclass[11pt,a4paper]{article}

\usepackage[utf8]{inputenc} % allow utf-8 input
\usepackage[T1]{fontenc}    % use 8-bit T1 fonts
\usepackage{hyperref}       % hyperlinks
\usepackage{url}            % simple URL typesetting
\usepackage{booktabs}       % professional-quality tables
\usepackage{amsfonts}       % blackboard math symbols
\usepackage{nicefrac}       % compact symbols for 1/2, etc.
\usepackage{microtype}      % microtypography
\usepackage{xcolor}         % colors
\usepackage{graphicx}       % graphics
\usepackage[normalem]{ulem} % strikethrough

\begin{document}

\section{Issue 1: Fraction as superscript base not grouped}

The equation below raises a fraction to the power 2.
Expected LaTeX: \{\textbackslash{}frac\{(x-c)\}\{v\}\}\textasciicircum{}2
Docling produces: \textbackslash{}frac\{(x-c)\}\{v\}\textasciicircum{}2

When <m:sSup> has an <m:f> (fraction) as its base <m:e>, docling emits
\textbackslash{}frac\{num\}\{den\}\textasciicircum{}2 without grouping braces. In LaTeX this is ambiguous:
some renderers apply \textasciicircum{}2 only to the last token, not the whole fraction.
The output should be \{\textbackslash{}frac\{(x-c)\}\{v\}\}\textasciicircum{}2 or \textbackslash{}left(\textbackslash{}frac\{(x-c)\}\{v\}\textbackslash{}right)\textasciicircum{}2.

$${\frac{(x-c)}{v}}^{2}$$

\end{document}